\documentclass[../../main.tex]{subfiles}
\begin{document}
\section{定向与覆盖映射}\label{section:ch15-coverings}

虽然构造一个定向通常很容易证明给定光滑流形可定向，但证明一个流形不可定向却要棘手得多。覆盖空间理论提供了做这件事最有用的技巧之一。本节中，我们探讨可定向性与覆盖映射之间的密切关系。

我们的第一个结果是拉回定向的一个简单应用。

\begin{proposition}\label{proposition:ch15-covering-orientable}
若 $\pi:E\to M$ 是一个光滑覆盖映射且 $M$ 可定向，则 $E$ 也可定向。
\end{proposition}

\begin{proof}
因为覆盖映射是局部微分同胚，这直接由命题\ref{proposition:ch15-pullback-orientation}得到。
\end{proof}

下一个定理更有趣。若 $G$ 是一个光滑作用在光滑流形 $E$ 上的 Lie 群（比如说在左侧作用），当对每个 $g\in G$，微分同胚 $x\mapsto g\cdot x$ 保定向时，就称这个作用是保定向作用。

\begin{theorem}\label{theorem:ch15-normal-covering}
设 $E$ 是一个连通的、有向的、带边或不带边的光滑流形，$\pi:E\to M$ 是一个光滑正规覆盖映射。则 $M$ 可定向当且仅当 $\operatorname{Aut}_\pi(E)$ 在 $E$ 上的作用是保定向的。
\end{theorem}

\begin{proof}
令 $\mathcal O_E$ 表示 $E$ 上给定的定向。首先设 $M$ 可定向，任取 $q\in E$。由于 $M$ 连通，它恰有两个定向，其中一个具有性质：$\mathrm{d}\pi_q:T_qE\to T_{\pi(q)}M$ 保定向。称那个定向为 $\mathcal O_M$。拉回定向 $\pi^*\mathcal O_M$ 在 $q$ 处与给定定向一致，故由命题\ref{proposition:ch15-two-orientations}它必等于 $\mathcal O_E$。设 $\varphi\in\operatorname{Aut}_\pi(E)$。$\pi\circ\varphi=\pi$ 这一事实蕴含
\begin{align*}
\varphi^*\mathcal O_E=\varphi^*(\pi^*\mathcal O_M)=(\pi\circ\varphi)^*\mathcal O_M=\pi^*\mathcal O_M=\mathcal O_E.
\end{align*}
因此 $\varphi$ 保定向。

反之，设 $\operatorname{Aut}_\pi(E)$ 的作用保定向，令 $p\in M$。若 $U\subseteq M$ 是 $p$ 的任意一个被均匀覆盖的邻域，则存在一个光滑截面 $\sigma:U\to E$，它在 $U$ 上诱导一个定向 $\sigma^*\mathcal O_E$。设 $\sigma_1:U\to E$ 是 $U$ 上任意另一个光滑局部截面。由于 $\pi$ 是正规覆盖，$\operatorname{Aut}_\pi(E)$ 在 $\pi$ 的每个纤维上传递地作用，故存在一个覆盖自同构 $\varphi$，使 $\sigma_1(p)=\varphi(\sigma(p))$。于是 $\varphi\circ\sigma$ 是 $\pi$ 的一个局部截面，它在 $p$ 处与 $\sigma_1$ 一致，因而在整个 $U$ 上 $\sigma_1=\varphi\circ\sigma$。由于 $\varphi$ 保定向，
\begin{align*}
\sigma_1^*\mathcal O_E=\sigma^*\varphi^*\mathcal O_E=\sigma^*\mathcal O_E,
\end{align*}
故由 $\sigma$ 与 $\sigma_1$ 诱导的定向相等。于是我们可以在 $M$ 上定义一个整体定向 $\mathcal O_M$：在每个被均匀覆盖的开子集上，把它定义为由任意局部截面诱导的拉回定向；上面的论证表明，这样定义的定向在重叠处一致。
\end{proof}

下面是上一定理的两个应用。

\begin{example}[射影空间的可定向性]\label{example:ch15-rp-orientable}
对 $n\geqslant1$，考虑例\ref{example:ch04-real-projective}中的光滑覆盖映射 $q:S^n\to\mathbb{RP}^n$。$q$ 唯一的非平凡覆盖自同构是对径映射 $\alpha(x)=-x$。习题\ref{exercise:ch15-03}表明 $\alpha$ 保定向当且仅当 $n$ 为奇数，故由此可知 $\mathbb{RP}^n$ 可定向当且仅当 $n$ 为奇数。
\end{example}

\begin{example}[Möbius 丛与 Möbius 带]\label{example:ch15-mobius-bundle}
设 $E$ 是 Möbius 丛（例\ref{example:ch10-03}）的全空间。用于定义 $E$ 的商映射 $q:\mathbb{R}^2\to E$ 是一个光滑正规覆盖映射，覆盖自同构群同构于 $\mathbb{Z}$，它在 $\mathbb{R}^2$ 上的作用为 $n\cdot(x,y)=(x+n,(-1)^n y)$。（你可以直接从定义验证这一点，也可以暂时接受它，等到我们在第~21 章发展出更多工具后再看一个更简单的证明；见习题\ref{exercise:ch21-09}。）当 $n$ 为奇数时，$\mathbb{R}^2$ 的微分同胚 $(x,y)\mapsto n\cdot(x,y)$ 把定向形式 $\mathrm{d}x\wedge\mathrm{d}y$ 拉回为 $-\mathrm{d}x\wedge\mathrm{d}y$，故 $\operatorname{Aut}_\pi(E)$ 的作用不保定向。于是定理\ref{theorem:ch15-normal-covering}表明 $E$ 不可定向。

对每个 $r>0$，矩形 $[0,1]\times[-r,r]$ 在 $q$ 下的像是一个 Möbius 带 $M_r$。由于 $q$ 限制为从 $\mathbb{R}\times[-r,r]$ 到 $M_r$ 的一个光滑覆盖映射，同样的论证表明 Möbius 带也不可定向。
\end{example}

\subsection{定向覆盖}

下面我们证明，每个不可定向光滑流形 $M$ 都有一个可定向的二叶覆盖流形。在点 $p\in M$ 上方的纤维将对应于 $T_pM$ 的两个定向。

为以统一方式处理可定向与不可定向两种情形，把我们对覆盖映射的定义稍作推广是有用的，即允许“覆盖空间”不连通。若 $N,M$ 是拓扑空间，设映射 $\pi:N\to M$ 满足覆盖映射的所有要求，只是 $N$ 可能不连通；这意味着 $N$ 局部道路连通，$\pi$ 满且连续，且每个 $p\in M$ 都有一个被 $\pi$ 均匀覆盖的邻域，就称 $\pi:N\to M$ 是一个\emph{广义覆盖映射}。若此外 $N,M$ 是带边或不带边的光滑流形且 $\pi$ 是局部微分同胚，就称它是一个\emph{广义光滑覆盖映射}。

\begin{lemma}\label{lemma:ch15-generalized-covering}
设 $N,M$ 是拓扑空间，$\pi:N\to M$ 是一个广义覆盖映射。若 $M$ 连通，则 $\pi$ 在 $N$ 的每个连通分支上的限制是一个覆盖映射。
\end{lemma}

\begin{proof}
设 $W$ 是 $N$ 的一个连通分支。若 $U$ 是 $M$ 的任意一个被 $\pi$ 均匀覆盖的开子集，则 $\pi^{-1}(U)$ 的每个连通分支都连通，因而含于 $N$ 的单个连通分支中。由此可知
\begin{align*}
(\pi|_W)^{-1}(U)=\pi^{-1}(U)\cap W
\end{align*}
或者是空集，或者是 $\pi^{-1}(U)$ 的一些连通分支的非空不交并，其中每个分支都被 $\pi|_W$ 同胚地映到 $U$ 上。特别地，这意味着 $\pi(W)$ 中每一点都有一个被 $\pi|_W$ 均匀覆盖的邻域。为完成证明，我们只需证 $\pi|_W$ 满。由于 $\pi$ 是局部同胚，$\pi(W)$ 是 $M$ 的一个开子集。另一方面，若 $p\in M\setminus\pi(W)$，$U$ 是 $p$ 的一个被 $\pi$ 均匀覆盖的邻域，则上一段的讨论表明 $(\pi|_W)^{-1}(U)=\varnothing$，这蕴含 $U\subseteq M\setminus\pi(W)$。因此 $\pi(W)$ 在 $M$ 中闭。由于 $W$ 非空，$\pi(W)=M$。
\end{proof}

设 $M$ 是一个连通的、光滑的、正维数的、带边或不带边的流形，令 $\widehat M$ 表示 $M$ 所有切空间的定向所成之集：
\begin{align*}
\widehat M=\{(p,\mathcal O_p):p\in M,\ \mathcal O_p\text{ 是 }T_pM\text{ 的一个定向}\}.
\end{align*}
把 $T_pM$ 的一个定向映到点 $p$ 本身，定义投影 $\widehat\pi:\widehat M\to M$：$\widehat\pi(p,\mathcal O_p)=p$。由于每个切空间恰有两个定向，这个映射的每条纤维的基数都是 $2$。映射 $\widehat\pi:\widehat M\to M$ 称为 $M$ 的\emph{定向覆盖}。

\begin{figure}[htbp]
\centering
\includegraphics[width=0.6\textwidth]{fig15-9-coverauto.png}
\caption{$\widehat M$ 的非平凡覆盖自同构}
\label{figure:ch15-coverauto}
\end{figure}

\begin{proposition}[定向覆盖的性质]\label{proposition:ch15-orientation-covering}
设 $M$ 是一个连通的、光滑的、正维数的、带边或不带边的流形，$\widehat\pi:\widehat M\to M$ 是它的定向覆盖。则 $\widehat M$ 可被赋予一个带边或不带边的、光滑的、有向流形的结构，具有下列性质：
\begin{enumerate}[(a)]
\item $\widehat\pi:\widehat M\to M$ 是一个广义光滑覆盖映射。
\item 一个连通开子集 $U\subseteq M$ 被 $\widehat\pi$ 均匀覆盖，当且仅当 $U$ 可定向。
\item 若 $U\subseteq M$ 是一个被均匀覆盖的开子集，则 $U$ 的每个定向都是 $\widehat\pi$ 在 $U$ 上的某个局部截面所诱导的拉回定向。
\end{enumerate}
\end{proposition}

\begin{proof}
我们首先通过为 $\widehat M$ 定义一组基来赋予它拓扑。对每个序对 $(U,\mathcal O)$（其中 $U$ 是 $M$ 的一个开子集，$\mathcal O$ 是 $U$ 上的一个定向），如下定义 $\widehat U_{\mathcal O}\subseteq\widehat M$：
\begin{align*}
\widehat U_{\mathcal O}=\{(p,\mathcal O_p)\in\widehat M:p\in U,\ \mathcal O_p\text{ 是 }\mathcal O\text{ 所确定的 }T_pM\text{ 的定向}\}.
\end{align*}
我们将证明，所有形如 $\widehat U_{\mathcal O}$ 的子集构成 $\widehat M$ 上一个拓扑的一组基。给定任意一点 $(p,\mathcal O_p)\in\widehat M$，令 $U$ 为 $p$ 在 $M$ 中的一个可定向邻域，$\mathcal O$ 是 $U$ 上的一个定向。必要时用 $-\mathcal O$ 代替 $\mathcal O$，可设给定定向 $\mathcal O_p$ 与 $\mathcal O$ 所确定的 $T_pM$ 的定向相同。由此可知 $(p,\mathcal O_p)\in\widehat U_{\mathcal O}$，故所有形如 $\widehat U_{\mathcal O}$ 的集族覆盖 $\widehat M$。若 $\widehat U_{\mathcal O}$ 与 $\widehat U'_{\mathcal O'}$ 是两个这样的集，$(p,\mathcal O_p)$ 是它们交集中的一点，则 $\mathcal O_p$ 是 $\mathcal O$ 与 $\mathcal O'$ 两者所确定的 $T_pM$ 的定向。若 $V$ 是 $U\cap U'$ 中含 $p$ 的连通分支，则限制后的定向 $\mathcal O|_V$ 与 $\mathcal O'|_V$ 在 $p$ 处一致，因而由命题\ref{proposition:ch15-two-orientations}相同，故 $\widehat U_{\mathcal O}\cap\widehat U'_{\mathcal O'}$ 包含基集 $\widehat V_{\mathcal O|_V}$。于是我们在 $\widehat M$ 上定义了一个拓扑。

注意，对每个可定向开子集 $U\subseteq M$ 与 $U$ 的每个定向 $\mathcal O$，$\widehat\pi$ 把基集 $\widehat U_{\mathcal O}$ 双射到 $U$。由于可定向开子集构成 $M$ 拓扑的一组基，这蕴含 $\widehat\pi$ 限制为从 $\widehat U_{\mathcal O}$ 到 $U$ 的同胚。特别地，$\widehat\pi$ 是局部同胚。

接着我们证明，在这个拓扑下 $\widehat\pi$ 是一个广义覆盖映射。设 $U\subseteq M$ 是一个可定向连通开子集，$\mathcal O$ 是 $U$ 的一个定向。则 $\widehat\pi^{-1}(U)$ 是开子集 $\widehat U_{\mathcal O}$ 与 $\widehat U_{-\mathcal O}$ 的不交并，且 $\widehat\pi$ 限制为从这两个集各自到 $U$ 的同胚。于是每个这样的集 $U$ 都被均匀覆盖，故 $\widehat\pi$ 是一个广义覆盖映射。由引理\ref{lemma:ch15-generalized-covering}，$\widehat\pi$ 限制为 $\widehat M$ 每个连通分支上的一个普通覆盖映射，故命题\ref{proposition:ch04-cover-manifold}表明每个这样的分支是一个带边或不带边的拓扑 $n$-流形，并具有唯一的光滑结构使 $\widehat\pi$ 成为一个光滑覆盖映射。这些光滑结构合起来给出整个 $\widehat M$ 上的一个光滑结构。这就完成了 (a) 的证明。

接着我们给 $\widehat M$ 一个定向。令 $\widehat p=(p,\mathcal O_p)$ 是 $\widehat M$ 中的一点。由定义，$\mathcal O_p$ 是 $T_pM$ 的一个定向；故我们可以给 $T_{\widehat p}\widehat M$ 唯一一个定向 $\widehat{\mathcal O}_{\widehat p}$，使 $\mathrm{d}\widehat\pi_{\widehat p}:T_{\widehat p}\widehat M\to T_pM$ 保定向。这就在 $\widehat M$ 上定义了一个逐点定向 $\widehat{\mathcal O}$。在每个基开子集 $\widehat U_{\mathcal O}$ 上，定向 $\widehat{\mathcal O}$ 与由 $\widehat\pi$（的限制）从 $(U,\mathcal O)$ 诱导的拉回定向一致，故它连续。

接着我们证 (b)。我们早先证明了 $M$ 的每个可定向连通开子集都被 $\widehat\pi$ 均匀覆盖。反之，若 $U\subseteq M$ 是任意一个被均匀覆盖的开子集，则由命题\ref{proposition:ch04-local-section}，存在 $\widehat\pi$ 的一个光滑局部截面 $\sigma:U\to\widehat M$，它由命题\ref{proposition:ch15-pullback-orientation}把 $\widehat{\mathcal O}$ 拉回到 $U$ 上的一个定向。

最后，为证 (c)，设 $U\subseteq M$ 被均匀覆盖，因而可定向。给定 $U$ 的任意定向 $\mathcal O$，通过令 $\sigma(p)=(p,\mathcal O_p)$ 定义一个截面 $\sigma:U\to\widehat M$。为看清 $\sigma$ 连续，设 $\widehat U'_{\mathcal O'}$ 是 $\widehat M$ 的任意基开子集。则对 $U\cap U'$ 的每个连通分支 $V$，限制后的定向 $\mathcal O|_V$ 与 $\mathcal O'|_V$ 必定要么在整个 $V$ 上一致，要么在整个 $V$ 上相反，故 $\sigma^{-1}(\widehat U'_{\mathcal O'})$ 是这些分支的一个并，因而是开集。
\end{proof}

\begin{theorem}[定向覆盖定理]\label{theorem:ch15-orientation-covering-theorem}
设 $M$ 是一个连通的、带边或不带边的光滑流形，$\widehat\pi:\widehat M\to M$ 是它的定向覆盖。
\begin{enumerate}[(a)]
\item 若 $M$ 可定向，则 $\widehat M$ 恰有两个连通分支，且 $\widehat\pi$ 在每个分支上的限制都是到 $M$ 的微分同胚。
\item 若 $M$ 不可定向，则 $\widehat M$ 连通，且 $\widehat\pi$ 是一个二叶光滑覆盖映射。
\end{enumerate}
\end{theorem}

\begin{proof}
若 $M$ 可定向，则命题\ref{proposition:ch15-orientation-covering}(b) 表明 $M$ 被 $\widehat\pi$ 均匀覆盖，这意味着 $\widehat M$ 有两个连通分支，每个都被微分同胚地映到 $M$ 上。

现在设 $M$ 不可定向。我们首先证明 $\widehat M$ 连通。令 $W$ 是 $\widehat M$ 的一个连通分支。引理\ref{lemma:ch15-generalized-covering}表明 $\widehat\pi|_W$ 是一个覆盖映射，故其纤维都有相同的基数。由于 $\widehat\pi$ 的纤维基数为 $2$ 且 $W$ 非空，$\widehat\pi|_W$ 的纤维基数必为 $1$ 或 $2$。若基数为 $1$，则 $\widehat\pi|_W$ 将是一个单射光滑覆盖映射，因而是一个微分同胚；其逆将是 $\widehat\pi$ 的一个光滑截面，这将在 $M$ 上诱导出一个定向。故基数必为 $2$，这蕴含 $W=\widehat M$。由于 $\widehat M$ 连通，由引理\ref{lemma:ch15-generalized-covering} $\widehat\pi$ 是一个覆盖映射；又由于它是局部微分同胚，它是一个光滑覆盖映射。
\end{proof}

定向覆盖有时称为 $M$ 的\emph{有向二重覆盖}。除我们这里给出的构造方式外，还有其他构造它的方法；但正如下一个定理所示，构造的具体细节无关紧要，因为它们都给出同构的覆盖流形。

\begin{theorem}[定向覆盖的唯一性]\label{theorem:ch15-covering-uniqueness}
设 $M$ 是一个不可定向的、连通的、带边或不带边的光滑流形，$\widehat\pi:\widehat M\to M$ 是它的定向覆盖。若 $\widetilde M$ 是一个带边或不带边的、有向的光滑流形，并容许一个二叶光滑覆盖映射 $\widetilde\pi:\widetilde M\to M$，则存在唯一的保定向微分同胚 $\varphi:\widetilde M\to\widehat M$，使 $\widehat\pi\circ\varphi=\widetilde\pi$。
\end{theorem}

\begin{proof}
见习题\ref{exercise:ch15-11}。
\end{proof}

借助稍多一点的覆盖空间理论，我们得到下面这个可定向性的充分拓扑条件。若 $G$ 是一个群，$H\subseteq G$ 是一个子群，$H$ 在 $G$ 中的\emph{指数}是 $H$ 在 $G$ 中的左陪集集的基数。（若 $H$ 是正规子群，它就是商群 $G/H$ 的基数。）

\begin{theorem}\label{theorem:ch15-fundamental-group}
设 $M$ 是一个连通的、带边或不带边的光滑流形，并设 $M$ 的基本群没有指数为 $2$ 的子群。则 $M$ 可定向。特别地，若 $M$ 单连通，则它可定向。
\end{theorem}

\begin{proof}
假设 $M$ 不可定向，令 $\widehat\pi:\widehat M\to M$ 为它的定向覆盖，在此情形下它是一个真正的覆盖映射。任取 $q\in\widehat M$，令 $p=\widehat\pi(q)\in M$。令 $\alpha:\widehat M\to\widehat M$ 为交换 $\widehat\pi$ 每条纤维中两个点的映射（图\ref{figure:ch15-coverauto}）。为证 $\alpha$ 光滑，设 $U\subseteq M$ 是任意一个被均匀覆盖的开子集，$U_0,U_1\subseteq\widehat M$ 是 $\widehat\pi^{-1}(U)$ 的两个连通分支。由于 $\widehat\pi$ 限制为从每个分支到 $U$ 的微分同胚，我们可写
\begin{align*}
\alpha|_{U_0}=(\widehat\pi|_{U_1})^{-1}\circ(\widehat\pi|_{U_0}),
\end{align*}
它光滑。类似地，$\alpha|_{U_1}$ 也光滑。由于所有这样的集 $U_0,U_1$ 构成 $\widehat M$ 的一个开覆盖，可知 $\alpha$ 光滑；又由于它满足 $\widehat\pi\circ\alpha=\widehat\pi$，它是一个覆盖自同构。事实上，由于一个覆盖自同构由它在一个点上的作用所唯一确定，$\alpha$ 就是自同构群 $\operatorname{Aut}_{\widehat\pi}(\widehat M)$ 中唯一的非平凡元素，因而该群等于二元群 $\{\operatorname{Id}_{\widehat M},\alpha\}$。由于该自同构群在纤维上传递地作用，$\widehat\pi$ 是一个正规覆盖映射。令 $H$ 表示 $\pi_1(M,p)$ 的子群 $\widehat\pi_*(\pi_1(\widehat M,q))$。覆盖空间理论的一个基本结果（例如见\cite{LeeTM}，第~12 章）是商群 $\pi_1(M,p)/H$ 同构于 $\operatorname{Aut}_{\widehat\pi}(\widehat M)$。因此 $H$ 在 $\pi_1(M,p)$ 中的指数为 $2$。
\end{proof}
\end{document}
